% !TeX program = pdflatex
\documentclass[11pt,a4paper]{amsart}

\newcommand{\notelanguage}{finnish}
% Shared preamble for course notes and exercise sets.

\usepackage[T1]{fontenc}
\usepackage{lmodern}
\providecommand{\notelanguage}{english}
\usepackage[finnish,english]{babel}
\AtBeginDocument{\expandafter\selectlanguage\expandafter{\notelanguage}}

\newcommand{\notesfinnishlanguage}{finnish}
\ifx\notelanguage\notesfinnishlanguage
\newcommand{\notestheoremname}{Lause}
\newcommand{\noteslemmaname}{Lemma}
\newcommand{\notespropositionname}{Propositio}
\newcommand{\notescorollaryname}{Seurauslause}
\newcommand{\notesdefinitionname}{Määritelmä}
\newcommand{\notesexamplename}{Esimerkki}
\newcommand{\notesproblemname}{Tehtävä}
\newcommand{\notesexercisename}{Tehtävä}
\newcommand{\notesremarkname}{Huomautus}
\newcommand{\notessolutionname}{Ratkaisu}
\else
\newcommand{\notestheoremname}{Theorem}
\newcommand{\noteslemmaname}{Lemma}
\newcommand{\notespropositionname}{Proposition}
\newcommand{\notescorollaryname}{Corollary}
\newcommand{\notesdefinitionname}{Definition}
\newcommand{\notesexamplename}{Example}
\newcommand{\notesproblemname}{Problem}
\newcommand{\notesexercisename}{Exercise}
\newcommand{\notesremarkname}{Remark}
\newcommand{\notessolutionname}{Solution}
\fi

\usepackage[a4paper,margin=30mm]{geometry}
\usepackage{microtype}
\usepackage{graphicx}
\usepackage{enumitem}

\newcommand{\exerciseheading}[2]{%
  \par\noindent
  {\large\bfseries #1\par}
  \vspace{0.1em}
  \noindent{\large #2\par}
  \vspace{1.5em}
}

\usepackage{amsmath}
\usepackage{amssymb}
\usepackage{amsthm}
\usepackage{mathtools}
\usepackage{aliascnt}

\usepackage[hidelinks,pdfusetitle]{hyperref}

\numberwithin{equation}{section}
\setlist{itemsep=0.25em,topsep=0.5em}

\theoremstyle{plain}
\newtheorem{theorem}{\notestheoremname}[section]
\newaliascnt{lemma}{theorem}
\newtheorem{lemma}[lemma]{\noteslemmaname}
\aliascntresetthe{lemma}
\newaliascnt{proposition}{theorem}
\newtheorem{proposition}[proposition]{\notespropositionname}
\aliascntresetthe{proposition}
\newaliascnt{corollary}{theorem}
\newtheorem{corollary}[corollary]{\notescorollaryname}
\aliascntresetthe{corollary}

\theoremstyle{definition}
\newaliascnt{definition}{theorem}
\newtheorem{definition}[definition]{\notesdefinitionname}
\aliascntresetthe{definition}
\newaliascnt{example}{theorem}
\newtheorem{example}[example]{\notesexamplename}
\aliascntresetthe{example}
\newaliascnt{problem}{theorem}
\newtheorem{problem}[problem]{\notesproblemname}
\aliascntresetthe{problem}
\newtheorem{exercise}{\notesexercisename}

\theoremstyle{remark}
\newaliascnt{remark}{theorem}
\newtheorem{remark}[remark]{\notesremarkname}
\aliascntresetthe{remark}

\newenvironment{solution}
{
\begin{proof}[\notessolutionname]}
  {
\end{proof}}

\usepackage[nameinlink,noabbrev]{cleveref}

\crefname{theorem}{theorem}{theorems}
\crefname{lemma}{lemma}{lemmas}
\crefname{proposition}{proposition}{propositions}
\crefname{corollary}{corollary}{corollaries}
\crefname{definition}{definition}{definitions}
\crefname{example}{example}{examples}
\crefname{problem}{problem}{problems}
\crefname{exercise}{exercise}{exercises}
\crefname{remark}{remark}{remarks}

\ifx\notelanguage\notesfinnishlanguage
\crefname{theorem}{lause}{lauseet}
\crefname{lemma}{lemma}{lemmat}
\crefname{proposition}{propositio}{propositiot}
\crefname{corollary}{seurauslause}{seurauslauseet}
\crefname{definition}{määritelmä}{määritelmät}
\crefname{example}{esimerkki}{esimerkit}
\crefname{problem}{tehtävä}{tehtävät}
\crefname{exercise}{tehtävä}{tehtävät}
\crefname{remark}{huomautus}{huomautukset}
\fi

\newcommand{\NN}{\mathbb{N}}
\newcommand{\NNone}{\mathbb{N}_1}
\newcommand{\ZZ}{\mathbb{Z}}
\newcommand{\QQ}{\mathbb{Q}}
\newcommand{\RR}{\mathbb{R}}
\newcommand{\CC}{\mathbb{C}}
\newcommand{\PP}{\mathbb{P}}

\DeclareMathOperator{\im}{im}
\DeclareMathOperator{\rank}{rank}
\DeclarePairedDelimiter{\abs}{\lvert}{\rvert}
\DeclarePairedDelimiter{\norm}{\lVert}{\rVert}
\newcommand{\set}[1]{\left\lbrace#1\right\rbrace}
\newcommand{\maxset}[1]{\max\set{#1}}
\newcommand{\minset}[1]{\min\set{#1}}
\newcommand{\powerset}[1]{\mathcal{P}\!\left(#1\right)}


\title{Analyysi I\\Harjoitus 6}
\author{}
\date{}

\begin{document}

\exerciseheading{Analyysi I}{Harjoitus 6}

\begin{exercise}
  Olkoon \(b \in \RR\). Määritellään funktio \(f\colon \RR \to \RR\)
  asettamalla \(f(x) = x^2 + x + 1\), kun \(x > 0\),
  \(f(x) = \abs{x - 2b}\), kun \(x < 0\), ja
  \(f(0) = 7 + b^2 + b^3\). Millä vakion \(b\) arvoilla funktiolla
  \(f\) on raja-arvo origossa?
\end{exercise}

\begin{solution}
  Olkoon \(b \in \RR\). Tutkitaan funktion \(f\) toispuoleiset
  raja-arvot origossa.

  Osoitetaan ensin, että \(\lim_{x \to 0^+} f(x) = 1\).
  Olkoon \(\varepsilon > 0\) ja valitaan
  \(\delta = \minset{1, \frac{\varepsilon}{2}} > 0\).
  Kun \(0 < x < \delta\), niin \(0 < x < 1\), joten \(x^2 < x\).
  Näin ollen
  \[
    \abs{f(x) - 1} = \abs{x^2 + x}
    \leq 2x < 2\delta \leq \varepsilon.
  \]
  Siis \(\lim_{x \to 0^+} f(x) = 1\).

  Osoitetaan seuraavaksi, että \(\lim_{x \to 0^-} f(x) = \abs{2b}\).
  Olkoon \(\varepsilon > 0\) ja valitaan \(\delta = \varepsilon > 0\).
  Kun \(-\delta < x < 0\), käänteisen kolmioepäyhtälön nojalla
  \[
    \begin{aligned}
      \abs{f(x) - \abs{2b}}
      &= \abs{\abs{x - 2b} - \abs{-2b}} \\
      &\leq \abs{(x - 2b) - (-2b)} \\
      &= \abs{x} < \delta = \varepsilon.
    \end{aligned}
  \]
  Siis \(\lim_{x \to 0^-} f(x) = \abs{2b}\).

  Lauseen 3.2.3 mukaan funktiolla \(f\) on raja-arvo origossa jos ja
  vain jos toispuoleiset raja-arvot ovat yhtä suuret. Edellä osoitetun
  perusteella tämä toteutuu jos ja vain jos \(\abs{2b} = 1\), eli
  \(b = \pm\frac{1}{2}\). Tällöin \(\lim_{x \to 0} f(x) = 1\).
  Funktion arvo \(f(0)\) ei vaikuta raja-arvon olemassaoloon eikä sen arvoon.
\end{solution}

\begin{exercise}
  Todista, että jos toinen raja-arvoista \(\lim_{x \to \infty} f(x)\)
  ja \(\lim_{x \to 0^+} f\left(\frac{1}{x}\right)\) on olemassa ja on
  reaaliluku, niin on toinenkin ja lisäksi raja-arvot ovat reaalilukuina
  yhtäsuuria.
\end{exercise}

\begin{exercise}
  Olkoon \(f\colon (1,2) \to \RR\) funktio, jolla on vasemmanpuoleinen
  raja-arvo pisteessä 2. Osoita, että on olemassa sellaiset
  \(m, \delta > 0\), että
  \[
    \abs{f(x)} < m
  \]
  kaikilla \(2 - \delta < x < 2\). Voidaanko aina valita, että \(m = 1\)?
\end{exercise}

\begin{exercise}
  Olkoon \(a\) positiivinen reaaliluku. Olkoon
  \(f\colon (0,\infty) \to \RR\), \(f(x) = \frac{a}{x}\).
  Osoita suoraan määritelmistä, että \(\lim_{x \to 0^+} f(x) = \infty\)
  ja \(\lim_{x \to \infty} f(x) = 0\).
\end{exercise}

\begin{exercise}
  Olkoon \(r \in \RR\), ja \(f\colon (a,b) \to \RR\) aidosti kasvava.
  Osoita, että \(rf\colon (a,b) \to \RR\) on aidosti kasvava jos
  \(r > 0\), aidosti vähenevä jos \(r < 0\), sekä kasvava ja vähenevä
  jos \(r = 0\).
\end{exercise}

\begin{exercise}
  Olkoon \(f\colon (a,b) \to \RR\) vähenevä, ja \(x_0 \in (a,b)\).
  Osoita, että
  \[
    \lim_{x \to x_0^-} f(x) \geq \lim_{x \to x_0^+} f(x).
  \]
  Voiko epäyhtälö olla aito?
\end{exercise}

\end{document}
